\section{Related works}
\label{related_works}
\paragraph{Cognitive map models.} 
Cognitive map models typically attribute structural abstraction to MEC, sensory binding to HPC, and sensory prediction to their interaction. 
The Tolman-Eichenbaum Machine (TEM)~\citep{TolmanEichenbaumMachineUnifying2020} extends beyond spatial domains by using recurrent networks to form cognitive maps that predict observations from predefined actions, though it requires re-learning abstract maps across environments. 
Clone-structured cognitive graphs (CSCG)~\citep{georgeClonestructuredGraphRepresentations2021} offer graph-based Markovian representations of structural relationships without prior constraints, but remain limited to discrete domains. 
Vector-HaSH~\citep{chandraEpisodicAssociativeMemory2025} generates velocity inputs from hippocampal states to drive grid cells forming episodic memories, but its velocity vectors are learned by memorizing the whole sequence.
All these works lack the critical step of inferring shared abstract structure from sequences in continuous and real-world environments.

\paragraph{Abstract velocity extraction.} In neuroscience, \cite{iyerflexible} employs an inverse model to extract abstract velocities from high-dimensional observations and maps them to low-dimensional grid cell velocity inputs. However, this approach significantly simplifies the complex representation learning in HPC-MEC circuits, focusing only on simple artificial stimuli.

\begin{figure*}[htbp]
    \vspace{-0.2 cm}
    \centering
    \includegraphics[width=1\textwidth]{Figures/methods/Model.png}
    \caption{\textbf{Overview of the model architecture.} (A) Video clips are passed through the visual encoder to obtain observation embeddings $\boldsymbol{s}$, which are encoded through the HPC to produce HPC embeddings $\boldsymbol{p}$, and then passed to the MEC to generate MEC embeddings $\boldsymbol{g}$. Finally, the generative pathway decodes them into the observation. The multi-scale VAE is fixed during training.
    (B) The latent transition $\boldsymbol{z}_t$ operates on the MEC embedding $\boldsymbol{g}_t$ at time $t$ using continuous attractor dynamics to generate the next MEC embedding $\boldsymbol{g}_{t+1}$.
    (C) The inverse model is used to infer latent transitions from the MEC embeddings $\boldsymbol{g}$.}
    \label{fig:methods:models}
    \vspace{-0.2 cm}
\end{figure*}

\paragraph{Infer latent transitions from the observations.}
World models~\citep{haWorldModels2018,lecunPathAutonomousMachine} are generative models that predict future observations based on past observations and actions. When ground-truth actions are unavailable, inverse and forward dynamics models are commonly employed to infer actions and predict future states from observation-only demonstrations~\citep{bruceGenieGenerativeInteractive2024, yeBECOMEPROFICIENTPLAYER2023, lapo}.
% However, the state-action interaction mechanisms in these models often lack interpretability. Common approaches include concatenating latent actions with corresponding frames~\citep{lapo, cuiDynaMoInDomainDynamics2024, yeBECOMEPROFICIENTPLAYER2023}, implementing cross-attention mechanisms~\citep{yeLatentActionPretraining2024, chenIGORImageGOalRepresentations2024}, or using additive embeddings~\citep{bruceGenieGenerativeInteractive2024, chenMotoLatentMotion2024}. Our approach implements state-action interaction through grid cell path integration, ensuring that latent actions function explicitly as velocity terms acting on current states.
Genie~\citep{bruceGenieGenerativeInteractive2024}, FICC~\citep{yeBECOMEPROFICIENTPLAYER2023}, and LAPO~\citep{lapo} primarily target 2D gaming environments, while
LAPA~\citep{yeLatentActionPretraining2024}, Moto~\citep{chenMotoLatentMotion2024}, IGOR~\citep{chenIGORImageGOalRepresentations2024}, and UniVLA~\citep{bu2025univla} focus on latent action extraction from real-world settings for pretraining policy models.
% Although LAPA, Moto, IGOR, and UniVLA extract latent actions using VQ-VAE objectives, they are not explicitly designed for interpretable state-action interaction and focus primarily on utilizing latent actions for control. Our approach implements state-action interaction through grid cell path integration, ensuring that latent transitions function explicitly as velocities acting on current states. 
AdaWorld~\citep{gaoAdaWorldLearningAdaptable2025b} achieves latent action transfer but relies on a large-scale pretrained video diffusion model. While their framework utilizes AdaLN~\citep{peebles2023scalable} for state-action interactions, our model leverages biologically-grounded CANN dynamics to achieve this integration. 
% Furthermore, we adopt a hierarchical architecture to emulate the functional HPC-MEC coupling mechanism.
% providing a more principled basis for structural abstraction.
